\SourceLeaf{323}%
\[
\begin{aligned}
b={}&\frac{r\cos\theta^3\cos U}{a\cos\varphi^2}\{1-[\tfrac1{12}+\frac{ee}{8\cos\varphi^2\cos\theta^2}(\cos P^2-\sin P^2\\
&\qquad-ee(5\cos P^2\sin P^2-\sin P^4))]qq+\tfrac1{12}ss+\tfrac18uu\}
\end{aligned}
\]
Für \(\cos\theta^3=(1-ee\sin P^2)^{\frac32}\) findet man leicht die vermittelst (4) so weit, wie hier nöthig ist, geführte Entwicklung
\[
\cos\theta^3=(1-ee\sin B^2)^{\frac32}(1-\frac{9e^4\cos P^2\sin P^2}{8\cos\varphi^2(1-ee\sin P^2)}\mathord{.}qq)
\]
wodurch die vorhergehende Gleichung sich verwandelt in
\[
\begin{aligned}
b={}&\frac{r(1-ee\sin B^2)^{\frac32}\cos U}{a\cos\varphi^2}\{1-[\tfrac1{12}+\frac{ee}{8\cos\varphi^2(1-ee\sin P^2)}(\cos P^2-\sin P^2\\
&\qquad+ee(4\cos P^2\sin P^2+\sin P^4))]qq+\tfrac1{12}ss+\tfrac18uu\}
\end{aligned}
\]
oder in einer etwas veränderten Form
\[
\begin{aligned}
b={}&\frac{rk^3\cos U}{a\cos\varphi^2}\{1-\frac{1}{24\cos\varphi^2(1-ee\sin P^2)}\\
&\qquad(2+ee-(8ee-14e^4)\sin P^2-9e^4\sin P^4)qq\\
&\qquad+\tfrac1{12}ss+\tfrac18uu\}
\end{aligned}
\]
Schreibt man nun noch hierin
\[
\begin{aligned}
&\frac{\cos\varphi\mathord{.}b}{\BareRoot{(1-ee\sin P^2)}}\quad\text{anstatt }q,\quad\text{wegen }(5)\\
&\frac{rr(1-ee\sin P^2)^2}{aa\cos\varphi^2}\quad\text{anstatt }ss,\quad\text{wegen }(8)\text{ und }(9)\\
&T\text{ und }t\text{ anstatt }U\text{ und }u,\quad\text{wegen }(6)\text{ und }(7)
\end{aligned}
\]
und zuletzt
\[
B\text{ für }P\text{ wegen }(4),
\]
was alles, ohne Nachtheil für die Genauigkeit geschehen kann, so erhält man (II)
\[
\begin{aligned}
b={}&\frac{rk^3}{a\cos\varphi^2}\mathord{.}\cos T\{1-\frac1{24k^4}(2+ee-(8ee-14e^4)\sin B^2-9e^4\sin B^4)\mathord{.}bb\\
&\qquad+\frac{k^4}{12aa\cos\varphi^2}\mathord{.}rr+\tfrac18tt\}
\end{aligned}
\]

\Article{25}
\begingroup\raggedright
\hspace*{1.2em}Aus den Gleichungen (1) und (3) erhellet, dass \(qqu\) von \(s^3\cos U^2\sin U\operatorname{tang}Q\), oder nach (11), von \(\frac{s^3\cos\varphi\cos U^2\sin U\mathord{.}\operatorname{tang}P}{\cos\theta}\) nur um eine Grösse fünfter Ordnung verschieden ist: es ist daher verstattet, die Gleichung (7) auch so zu schreiben\par\endgroup
\[
t=u(1-\frac{ee\cos P^2}{6\cos\varphi^2}\mathord{.}qq)
\]
